%% Section VI draft: Conclusion.
%% Every number quoted here appears in Tables III-VI; nothing new is claimed.

\section{Conclusion}\label{sec:conclusion}
This paper posed cold-start magnetic-anomaly navigation as a single estimation
problem: an aircraft with an uncalibrated magnetometer installation must
recover its position from a scalar reading in which the map signal and the
aircraft's own field are mixed, with no calibration flight and no satellite
aiding. The formulation carries the Tolles--Lawson compensation coefficients as
time-varying variables of a factor graph alongside the inertial error state,
ties the two chains through the map measurement, and solves the graph with an
incremental fixed-lag smoother built on iSAM2, at a \SI{300}{\second} lag and a
\SI{1}{\hertz} state cadence. The estimator produces two outputs, a causal
estimate with no look-ahead and a committed estimate revised by the full lag,
and we report both throughout, because they answer different questions: the
causal output is what a navigator reads in flight, the committed output is what
a system that may lag real time delivers.

On the SGL~2020 flight data the answers are as follows. Against a recursive EKF
carrying the identical models, the causal output alone is more accurate on six
of eight cold-start cases, so windowed relinearization pays before any
look-ahead is spent. Against the strongest published causal method, an EKF
augmented with an online neural network, the linear-basis smoother is more
accurate on both magnetometers of the representative line with no learned
component at all. The committed output is the most accurate estimate on all
eight counted cases, between \SI{10.5}{} and \SI{28.4}{\meter} where
free-inertial drift is hundreds of metres, at \num{1.5} to \num{5.2} times the
accuracy of the best causal baseline. The lag sweep locates the design knee
near \SI{300}{\second}, about \SI{20}{\kilo\meter} of track, and shows that the
causal output is indifferent to the lag, so the choice of lag is purely a
purchase of smoothing accuracy with latency. The ablation finds no delicate
setting: one configuration serves every line and both sensors, and the update
costs \SIrange{6.0}{9.0}{\milli\second} per state in an unoptimized
implementation, two orders of magnitude inside real time.

Several limitations bound the study. Each line is a single physical
realization, so every DRMS is a point estimate; the statistical statements are
counts across cases, not distributions. The reported covariance of the
incremental smoother has not been validated for consistency on flight data,
and we do not report uncertainty bounds. The
baselines are our own implementations under one shared protocol; the
particle-filter reference in particular is reported as unstable under per-case
tuning, and we do not claim it represents the best achievable particle filter.
None of the lines is a long, straight-and-level operational transit, so
performance under the weak regressor excitation of a pure cruise is not
established. And the compensation prior that a cold start needs is a measured
setting rather than a derived one: we show a broad insensitive range, but give
no rule for choosing its scale from installation data.

These limitations set the agenda: consistency evaluation of the reported
covariance, a transit-like flight regime, sparser instantiation of the
compensation chain in the manner of~\cite{park2024}, and a native-factor
implementation that would carry the real-time margin from the state cadence
down to the raw sensor rate.
